%% 毕业设计小结

\begin{designsummary}
    这段围绕代码、数据和实验结果不断推进的毕业设计经历，如今已接近收尾。回顾整个过程，从最初理解3D Gaussian Splatting的基本原理，到逐步梳理项目代码、设计模块进行实验、组织训练与评估流程，再到最终完成论文撰写，每一步都伴随着查阅资料、调试程序和反复修改。这个过程不仅是一次毕业设计任务的完成，也是对四年专业学习成果的一次集中检验。
    
    通过本次毕业设计，我对三维重建与新视角合成任务有了更加系统的认识。过去在课堂中学习的计算机视觉、深度学习、和程序设计等知识，在课题实现中变成了具体的问题：如何组织多视角数据，如何利用COLMAP位姿结果初始化场景，如何在3DGS训练中加入FAME远区域增强，如何设计有效秩与体积正则，如何分析新视角合成质量。将这些问题逐步拆解并落实到代码和实验中，使我对科研问题和工程实现之间的关系有了更直接的体会。
    
    当然，本次毕业设计仍然存在不足。受时间和个人能力限制，系统在几何重建质量、复杂光照鲁棒性等方面仍有进一步优化空间；部分实验分析也还可以结合更多场景和更充分的可视化结果进行展开。这些尚存的不足，如同前路里待探寻的坐标，让我明晰了日后精进的路径。
    
    总体而言，本次毕业设计让我在专业知识、工程实践和科研表达三个方面都有了明显收获。它既是本科阶段学习成果的总结，也为我今后继续学习计算机视觉相关技术打下了基础。未来我将继续保持对技术问题的探索热情，在实际工作和进一步学习中不断完善自己的知识结构和实践能力。
    \end{designsummary}
    